\documentclass[11pt]{article}
\usepackage[T1]{fontenc}
\usepackage[utf8]{inputenc}
\usepackage{lmodern}
\usepackage[margin=1.1in]{geometry}
\usepackage{amsmath,amssymb,amsthm,mathtools}
\usepackage{microtype,xcolor,titlesec,hyperref}

\titleformat{\section}{\normalfont\LARGE\bfseries}{\thesection}{1em}{}
\titleformat{\subsection}{\normalfont\Large\bfseries}{\thesubsection}{1em}{}

\hypersetup{
  pdftitle={Generic Degree Two in JC(2): Secant Projectors and Galois Rigidity},
  pdfauthor={Chloe van der Vlugt},
  colorlinks=true,
  linkcolor=blue!55!black,
  citecolor=green!40!black,
  urlcolor=blue!60!black
}

\newtheorem{theorem}{Theorem}[section]
\newtheorem{proposition}[theorem]{Proposition}
\newtheorem{corollary}[theorem]{Corollary}

\DeclareMathOperator{\Spec}{Spec}\DeclareMathOperator{\Frac}{Frac}
\DeclareMathOperator{\Gal}{Gal}

\newcommand{\A}{\mathbb A}
\newcommand{\C}{\mathbb C}
\newcommand{\Obs}{\operatorname{Obs}}
\newcommand{\leanpath}[1]{\texttt{\nolinkurl{#1}}}

\newenvironment{leanfiles}
  {\par\medskip\begingroup\footnotesize\raggedright\noindent
   \textit{Lean correspondence.}\ }
  {\par\endgroup\medskip}

\title{\texorpdfstring
  {Generic Degree Two in \(JC(2)\):\\
   Secant Projectors and Galois Rigidity}
  {Generic Degree Two in JC(2): Secant Projectors and Galois Rigidity}}
\author{Chloe van der Vlugt}\date{\today}

\begin{document}

\maketitle

\begin{abstract}
We study planar polynomial maps of generic degree two through the scheme of
pairs with the same image.  For
\[
  F=(P,Q)\colon\A^2_{\C}\longrightarrow\A^2_{\C},
\]
put \(A=\C[x_1,x_2]\), let \(I_R(F)\) be the collision ideal in the
pair ring \(S=\C[x_1,x_2,y_1,y_2]\), and let \(I_\Delta\) be the diagonal
ideal.  Set \(C_F=S/I_R(F)\).  Their quotient
\[
  \Obs(F)=I_\Delta/I_R(F)\subseteq C_F
\]
is the ideal of the diagonal in the collision ring.  When \(F\) is Keller,
an ordered divided-difference matrix produces an explicit idempotent
\(p_F\in C_F\) with \(\Obs(F)=C_Fp_F\).

Writing \(B=\C[P,Q]\subseteq A\), \(K=\Frac(B)\), and
\(L=\Frac(A)\), suppose that the induced extension \(L/K\) has degree
two.  Then the generic collision algebra decomposes compatibly with the
diagonal as
\[
  K\otimes_BC_F\simeq_K L\times L,
\]
where diagonal evaluation is first projection.  The second factor forces
\(\Obs(F)\neq0\); when \(F\) is Keller, the secant identity
\(\Obs(F)=C_Fp_F\) then gives \(p_F\neq0\).  On the other hand, a separable
quadratic extension is Galois, and classical Keller--Galois rigidity makes
a hypothetical generic-degree-two Keller map a polynomial automorphism.  Then
\(\Obs(F)=0\) and \(p_F=0\), a contradiction.  Thus no planar Keller map
has generic degree two.  The collision and secant arguments are formalized
in Lean; the classical rigidity theorem is isolated there as the sole
external mathematical input to the degree-two exclusion.
\end{abstract}

\begin{center}
  \small Licensed under
  \href{https://creativecommons.org/licenses/by/4.0/}
    {CC BY 4.0}.
\end{center}

\tableofcontents

\section{Introduction}
\label{sec:introduction}

The complex two-dimensional Jacobian conjecture \(JC(2)\) asks whether every
polynomial map
\[
  F=(P,Q)\colon\A^2_{\C}\longrightarrow\A^2_{\C},
  \qquad \det JF\in\C^\times,
\]
has a polynomial inverse.  Such an \(F\) is called a \emph{Keller map}.
The condition on the Jacobian is local, while polynomial invertibility is
global.  Collision geometry makes the global issue visible by retaining
pairs \((x,y)\) for which \(F(x)=F(y)\).

The companion paper \cite{VanDerVlugt2026CubicCollision} develops the
dimension-independent collision construction.  It associates to a
polynomial map two ideals
\[
  I_R(F)=\bigl(F_i(x)-F_i(y)\bigr),
  \qquad
  I_\Delta=\bigl(x_i-y_i\bigr)
\]
inside the coordinate ring of ordered pairs.  The containment
\(I_R(F)\subseteq I_\Delta\) induces the obstruction ideal
\[
  \Obs(F)=I_\Delta/I_R(F)
\]
in the collision ring \(C_F\).  This ideal defines the diagonal inside the
collision scheme.  Its vanishing says that the collision scheme consists
only of the diagonal; over \(\C\), this is equivalent to polynomial
automorphy.

The same construction has a generic field-theoretic form.  Write
\[
  A=\C[x_1,x_2],\qquad
  B=\C[P,Q]\subseteq A,\qquad
  K=\Frac(B),\qquad
  L=\Frac(A).
\]
For a Keller map, \(L/K\) is finite and separable, and
\[
  d_F=[L:K]
\]
is the generic degree of \(F\).  At the generic point, Paper~I gives the
compatible isomorphism
\[
  K\otimes_BC_F\simeq_KL\otimes_KL,
\]
with base-changed diagonal evaluation equal to tensor multiplication.

Paper~I already uses classical Galois rigidity to exclude generic degrees
one and two in every dimension.  The purpose here is to realize the
quadratic exclusion explicitly in the planar collision ring.  A recent
preprint states the stronger exclusion of every prime generic degree
\cite{Moskowicz2024PrimeDegree}.  We do not use that broader argument; the
focus here is the explicit collision certificate and its formal verification.

This paper specializes that construction to \(d_F=2\).  There are two
complementary descriptions of the collision obstruction.  The first is
geometric and explicit: ordered divided differences produce the planar
multivariate B\'ezoutian determinant and an idempotent \(p_F\) generating
\(\Obs(F)\).  The second is generic
and field-theoretic: a separable quadratic extension satisfies
\[
  L\otimes_KL\simeq_L L\times L,
\]
and tensor multiplication is projection onto the marked first factor.
For a Keller map, the second factor is a genuine residual generic
collision, forcing \(p_F\neq0\).

Classical Galois rigidity gives the opposite conclusion.  Since every
separable quadratic extension is normal, a hypothetical planar Keller map
with \(d_F=2\) would be a polynomial automorphism
\cite{Campbell1973,Razar1979,Wright1981}.  Its collision scheme would then
be the diagonal, so \(p_F=0\).  The contradiction excludes generic degree
two.

Section~\ref{sec:collision-inputs} recalls, without reproving, the part of
the general construction needed here.  Section~\ref{sec:secant-projector}
constructs the planar secant projector.  Section
\ref{sec:quadratic-collision} proves quadratic collision nonvanishing, and
Section~\ref{sec:rigidity-contradiction} combines it with Galois rigidity.
Section~\ref{sec:lean-status} records the formalization boundary.

This is a generic-degree-two theorem, not a proof of all of \(JC(2)\).
The remaining possible counterexamples have nonnormal function-field
extensions of generic degree at least three.  No normalization, boundary,
or inertia construction is needed below.

\section{Inputs from Paper I}
\label{sec:collision-inputs}

We use the dimension-independent construction of Paper~I
\cite{VanDerVlugt2026CubicCollision}.  To keep the two manuscripts
notationally identical, let
\[
  F=(F_1,\dots,F_n)\colon\A^n_{\C}\longrightarrow\A^n_{\C},
\]
and put
\[
  \begin{gathered}
  A=\C[x_1,\dots,x_n],\qquad
  S=\C[x_1,\dots,x_n,y_1,\dots,y_n],\\
  I_R(F)=\bigl(F_i(x)-F_i(y)\bigr)_{i=1}^n,\qquad
  I_\Delta=\bigl(x_i-y_i\bigr)_{i=1}^n,\\
  C_F=S/I_R(F),\qquad
  R_F=\Spec C_F.
  \end{gathered}
\]
Since \(I_R(F)\subseteq I_\Delta\), diagonal substitution descends along
the quotient map \(q_F\colon S\twoheadrightarrow C_F\) to
\[
  \bar\mu_F\colon C_F\twoheadrightarrow A,
  \qquad [s]\longmapsto s(x,x).
\]
The image of \(I_\Delta\) is the ideal
\[
  \Obs(F):=q_F(I_\Delta)=I_\Delta/I_R(F)\subseteq C_F,
\]
which defines the diagonal inside \(R_F\).

The following two results are quoted from Paper~I; no proof of the general
collision construction is repeated here.

\begin{theorem}[Diagonal kernel and automorphism detection]
\label{thm:quoted-diagonal-kernel}
For every complex polynomial map as above,
\[
  \ker(\bar\mu_F)=\Obs(F)=I_\Delta/I_R(F).
\]
Moreover,
\[
  F\text{ is a polynomial automorphism}
  \quad\Longleftrightarrow\quad
  I_R(F)=I_\Delta
  \quad\Longleftrightarrow\quad
  \Obs(F)=0.
\]
\end{theorem}

The implication \(\Obs(F)=0\Rightarrow F\) is a polynomial automorphism
uses Ax--Grothendieck
\cite{Ax1969}; below we need only the elementary forward implication.

Let
\[
  B=\C[F_1,\dots,F_n]\subseteq A,\qquad
  K=\Frac(B),\qquad
  L=\Frac(A),
\]
and let
\[
  \mu_L\colon L\otimes_KL\longrightarrow L,
  \qquad \ell\otimes\ell'\longmapsto\ell\ell'
\]
be tensor multiplication.

\begin{theorem}[Generic collision bridge]
\label{thm:quoted-generic-bridge}
If \(L/K\) is finite, then
\[
  K\otimes_BA\simeq_KL
\]
and there is a \(K\)-algebra isomorphism
\[
  \Phi_F\colon K\otimes_BC_F\xrightarrow{\ \sim\ }L\otimes_KL
\]
which carries the base-changed diagonal evaluation to \(\mu_L\).
Furthermore, flatness of the localization \(B\to K\) identifies
\[
  K\otimes_B\Obs(F)
  \simeq
  \ker(1\otimes\bar\mu_F)
\]
as \(K\)-vector spaces.
\end{theorem}

These statements combine the fiber-product, generic-base-change, and
diagonal-kernel results of Paper~I.  Their compatibility with \(\mu_L\)
marks which generic factor is the diagonal.

Return now to
\[
  F=(P,Q)\colon\A^2_{\C}\longrightarrow\A^2_{\C}.
\]
For a Keller map, \(P,Q\) are algebraically independent and \(L/K\) is
finite separable.  We write
\[
  d_F:=[L:K]
\]
for its generic degree.  The \(2\) in \(JC(2)\) denotes ambient dimension;
the separate condition \(d_F=2\) concerns the generic fiber.  The
nonvanishing argument below only uses a finite separable extension of
degree two.  The Keller condition enters through the nonzero constant
Jacobian and, at the end, through Galois rigidity.

\section{The explicit planar secant projector}
\label{sec:secant-projector}

The matrix below is the ordered multivariate B\'ezoutian matrix associated
to \(F\); compare \cite{McKean2023Bezoutians,SchejaStorch1975}.  We retain
the term \emph{secant matrix} because its determinant separates the
diagonal collision sheet from its complement.

Let
\[
  F=(P,Q)\colon\A^2_{\C}\longrightarrow\A^2_{\C},
  \qquad
  \det JF=c\in\C^\times.
\]
Set
\[
  S=\C[x_1,x_2,y_1,y_2],\qquad
  \mathbf d=
  \begin{pmatrix}x_1-y_1\\x_2-y_2\end{pmatrix},\qquad
  \mathbf r=
  \begin{pmatrix}P(x)-P(y)\\Q(x)-Q(y)\end{pmatrix}.
\]

\subsection{Ordered divided differences}

For \(H\in\C[x_1,x_2]\), define
\[
  D_1H
  =
  \frac{H(x_1,x_2)-H(y_1,x_2)}{x_1-y_1},
\]
\[
  D_2H
  =
  \frac{H(y_1,x_2)-H(y_1,y_2)}{x_2-y_2}.
\]
Both quotients lie in \(S\).  The ordered telescoping identity is
\[
  H(x)-H(y)
  =D_1H\,(x_1-y_1)+D_2H\,(x_2-y_2).
\]
Define
\[
  M_F=
  \begin{pmatrix}
    D_1P&D_2P\\
    D_1Q&D_2Q
  \end{pmatrix}.
\]
Then
\[
  \mathbf r=M_F\mathbf d,\qquad
  M_F\big|_\Delta=JF.
\]
Let
\[
  \delta_F:=\det(M_F)\in S.
\]

\begin{proposition}[Secant annihilation]
\label{prop:secant-annihilation}
The secant determinant satisfies
\[
  \delta_FI_\Delta\subseteq I_R(F),
  \qquad
  \delta_F\big|_\Delta=c.
\]
\end{proposition}

\begin{proof}
Multiplying \(\mathbf r=M_F\mathbf d\) by the adjugate matrix gives
\[
  \delta_F\mathbf d=\operatorname{adj}(M_F)\mathbf r.
\]
Both entries on the right belong to \(I_R(F)\).  Since the entries of
\(\mathbf d\) generate \(I_\Delta\), the first assertion follows.
Restriction to the diagonal gives
\[
  \delta_F\big|_\Delta
  =\det(M_F|_\Delta)
  =\det JF
  =c.
\]
\end{proof}

\subsection{The idempotent}

Let \(\bar\delta_F:=q_F(\delta_F)\) be the class of \(\delta_F\) in
\(C_F\), and put
\[
  e_F:=c^{-1}\bar\delta_F,\qquad
  p_F:=1-e_F=1-c^{-1}\bar\delta_F.
\]
The element \(e_F\) is the diagonal projector, while \(p_F\) is the
off-diagonal collision projector.

\begin{theorem}[Planar secant projector]
\label{thm:planar-secant-projector}
With the notation above,
\[
  e_F^2=e_F,\qquad
  p_F^2=p_F,\qquad
  e_Fp_F=0,\qquad
  e_F+p_F=1.
\]
Moreover,
\[
  \Obs(F)=C_Fp_F,\qquad
  \operatorname{Ann}_{C_F}(\Obs(F))=C_Fe_F,
\]
and \(p_F\) is the unique idempotent \(u\in C_F\) satisfying
\(\Obs(F)=C_Fu\).  Hence
\[
  C_F\simeq A\times C_Fp_F.
\]
In particular,
\[
  p_F=0\quad\Longleftrightarrow\quad\Obs(F)=0.
\]
\end{theorem}

\begin{proof}
Proposition~\ref{prop:secant-annihilation} gives
\[
  \bar\delta_F\Obs(F)=0,\qquad
  \bar\delta_F-c\in\Obs(F).
\]
Equivalently,
\[
  e_F\Obs(F)=0,\qquad e_F-1\in\Obs(F).
\]
Hence \(e_F(e_F-1)=0\), so \(e_F^2=e_F\); the corresponding identities
for \(p_F=1-e_F\) follow.

The second relation says \(p_F\in\Obs(F)\).  If
\(j\in\Obs(F)\), then
\[
  j=(e_F+p_F)j=p_Fj,
\]
so \(\Obs(F)=C_Fp_F\).  An element annihilates \(C_Fp_F\) exactly when its
\(p_F\)-part is zero, giving
\[
  \operatorname{Ann}_{C_F}(\Obs(F))=C_Fe_F.
\]

For uniqueness, let \(u^2=u\) and suppose \(C_Fu=C_Fp_F\).  Since
\(u\in C_Fp_F\), one has \(up_F=u\); since \(p_F\in C_Fu\), one has
\(up_F=p_F\).  Thus \(u=p_F\).

Complementary idempotents give
\[
  C_F\xrightarrow{\ \sim\ }C_Fe_F\times C_Fp_F,\qquad
  a\longmapsto(ae_F,ap_F).
\]
Diagonal evaluation restricts to an isomorphism \(C_Fe_F\simeq A\):
if \(a=\bar\mu_F(r)\), then \(a=\bar\mu_F(re_F)\), while its kernel is
\(C_Fe_F\cap\Obs(F)=0\).  Thus \(C_F\simeq A\times C_Fp_F\).
Finally, \(\Obs(F)=C_Fp_F\) gives
\[
  p_F=0\Longleftrightarrow\Obs(F)=0.
\]
\end{proof}

Abstractly, the existence of an idempotent splitting is expected: a Keller
map is unramified, so its diagonal in the collision scheme is open, while a
morphism between affine schemes has closed diagonal
\cite{StacksUnramifiedDiagonalOpen,StacksAffineDiagonalClosed}.  The theorem
supplies the explicit B\'ezoutian representative and identifies the ideal it
generates.  The uniqueness just proved shows that the resulting \(p_F\) is
independent of the ordering used to write the divided differences.

\section{Quadratic generic collisions}
\label{sec:quadratic-collision}

We now work at the generic point.  The essential field-theoretic fact is
the degree-two marked-root decomposition.

\subsection{The quadratic self-tensor product}

Let \(L/K\) be finite separable of degree two.  Paper~I's marked-root
tensor decomposition applies to a primitive generator \(\alpha\): it
splits the marked root from the residual polynomial
\[
  m_\alpha(T)=(T-\alpha)h_\alpha(T)
  \qquad\text{in }L[T].
\]
Here \(h_\alpha\) has degree one, so its quotient algebra is another copy
of \(L\).  Thus the general result of Paper~I
\cite{VanDerVlugt2026CubicCollision} specializes as follows.

\begin{proposition}[Quadratic marked-root decomposition]
\label{prop:quadratic-tensor-product}
There is an \(L\)-algebra isomorphism
\[
  \Psi_\alpha\colon
  L\otimes_KL\xrightarrow{\ \sim\ }L\times L
\]
under which tensor multiplication
\[
  \mu_L\colon L\otimes_KL\longrightarrow L
\]
is first projection.  Consequently,
\[
  \Psi_\alpha\bigl(\ker(\mu_L)\bigr)=0\times L\neq0.
\]
\end{proposition}

\subsection{Nonvanishing of the affine obstruction}

Combining the quoted generic collision bridge with
Proposition~\ref{prop:quadratic-tensor-product} transports the nonzero
second factor back to the affine obstruction.

\begin{theorem}[Quadratic collision nonvanishing]
\label{thm:quadratic-obstruction-nonzero}
Let \(F\colon\A^n_{\C}\to\A^n_{\C}\) be a polynomial map whose induced
extension \(L/K\) is finite separable of degree two.  Then
\[
  \Obs(F)\neq0.
\]
Equivalently,
\[
  I_R(F)\subsetneq I_\Delta.
\]
\end{theorem}

\begin{proof}
Flatness of \(B\to K\) gives a canonical identification
\[
  K\otimes_B\Obs(F)
  \simeq
  \ker(1\otimes\bar\mu_F).
\]
Under the compatible quadratic product, the right-hand side is
\[
  \ker\bigl(L\times L\xrightarrow{\operatorname{pr}_1}L\bigr)
  =0\times L,
\]
which is nonzero.  Therefore \(K\otimes_B\Obs(F)\neq0\), and hence
\(\Obs(F)\neq0\).  Since \(\Obs(F)=I_\Delta/I_R(F)\), this is equivalent
to strict containment.
\end{proof}

This theorem is independent of the Keller condition: generic degree two
already forces a residual generic collision.  For a planar Keller map, the
secant projector makes that factor explicit.

\begin{corollary}[Nonzero planar projector in generic degree two]
\label{cor:quadratic-projector-nonzero}
Let
\[
  F=(P,Q)\colon\A^2_{\C}\longrightarrow\A^2_{\C},
  \qquad
  \det JF=c\in\C^\times.
\]
If \(d_F=2\), then
\[
  p_F\neq0.
\]
\end{corollary}

\begin{proof}
The Keller condition makes \(L/K\) finite separable.
Theorem~\ref{thm:quadratic-obstruction-nonzero} gives
\(\Obs(F)\neq0\), and
Theorem~\ref{thm:planar-secant-projector} gives
\(\Obs(F)=C_Fp_F\).  Therefore \(p_F\neq0\).
\end{proof}

\section{Galois rigidity and the collision contradiction}
\label{sec:rigidity-contradiction}

We now use the classical Galois case of the Jacobian conjecture.

\begin{theorem}[Classical Keller--Galois rigidity]
\label{thm:keller-galois-rigidity}
Let
\[
  F\colon\A^2_{\C}\longrightarrow\A^2_{\C}
\]
be a Keller map.  If the induced finite separable extension \(L/K\) is
normal, equivalently Galois, then \(F\) is a polynomial automorphism.
\end{theorem}

This is the classical normal/Galois case
\cite{Campbell1973,Razar1979,Wright1981}; we use it as a literature theorem
and do not reprove it here.

Every separable extension of degree two is normal: if the minimal polynomial
of a primitive generator has one root in \(L\), its other root also lies in
\(L\).  Thus Theorem~\ref{thm:keller-galois-rigidity} applies automatically
on the generic-degree-two Keller locus.

\begin{proposition}[Rigidity forces collision vanishing]
\label{prop:rigidity-projector-zero}
Let
\[
  F=(P,Q)\colon\A^2_{\C}\longrightarrow\A^2_{\C},
  \qquad
  \det JF=c\in\C^\times.
\]
If \(d_F=2\), then Keller--Galois rigidity gives
\[
  \Obs(F)=0,\qquad p_F=0.
\]
\end{proposition}

\begin{proof}
The extension \(L/K\) is finite separable of degree two, hence Galois.
Theorem~\ref{thm:keller-galois-rigidity} makes \(F\) a polynomial
automorphism.  Theorem~\ref{thm:quoted-diagonal-kernel} then gives
\[
  I_R(F)=I_\Delta,\qquad \Obs(F)=0.
\]
Finally, Theorem~\ref{thm:planar-secant-projector} gives \(p_F=0\).
\end{proof}

\begin{theorem}[Exclusion of generic degree two]
\label{thm:no-planar-keller-degree-two}
No planar Keller map has generic degree two.  Equivalently, every
polynomial map
\[
  F=(P,Q)\colon\A^2_{\C}\longrightarrow\A^2_{\C},
  \qquad
  \det JF\in\C^\times,
\]
satisfies
\[
  d_F\neq2.
\]
\end{theorem}

\begin{proof}
Suppose that \(d_F=2\).
Corollary~\ref{cor:quadratic-projector-nonzero}, which reads the residual
factor in
\[
  K\otimes_BC_F\simeq_K L\times L,
\]
gives \(p_F\neq0\).  Proposition~\ref{prop:rigidity-projector-zero} gives
\(p_F=0\).  This is a contradiction.
\end{proof}

Equivalently, the two halves of the proof say
\[
  d_F=2\Longrightarrow\Obs(F)\neq0
\]
by the quadratic generic collision product, while
\[
  \text{\(F\) Keller and \(L/K\) normal}
  \Longrightarrow\Obs(F)=0
\]
by Galois rigidity and automorphism detection.  The secant construction
expresses both statements as nonvanishing and vanishing of one canonical
idempotent.

\section{Lean formalization status}
\label{sec:lean-status}

\begin{leanfiles}
The focused theorem spine is formalized in Lean~4.28.0 and Mathlib~v4.28.0
\cite{Lean4,Mathlib}.  The general collision objects are defined in
\leanpath{CollisionIdeals/General/Collision}.  The function-field
localization and generic collision bridge are in
\leanpath{CollisionIdeals/General/GenericFiber/FunctionField.lean}
and
\leanpath{CollisionIdeals/General/GenericFiber/CollisionDiagonal.lean}.

The degree-two marked-root calculation is proved in
\leanpath{CollisionIdeals/General/GenericFiber/QuadraticDecomposition.lean}.
Its composition with the polynomial generic-collision bridge, including
\[
  K\otimes_BC_F\simeq_K L\times L
\]
and the descent to \(\Obs(F)\neq0\), is in
\leanpath{CollisionIdeals/General/GenericFiber/QuadraticCollision.lean}.

The coefficientwise divided differences, secant determinant, annihilation
relation, diagonal restriction, and canonical projector are proved in
\leanpath{CollisionIdeals/Planar/ExplicitSecant.lean}.  The planar
nonvanishing and vanishing statements are in
\leanpath{CollisionIdeals/Planar/GenericDegreeTwo/Nonvanishing.lean}
and
\leanpath{CollisionIdeals/Planar/GenericDegreeTwo/Vanishing.lean}.
Their collision-native contradiction and \(d_F\neq2\) are assembled in
\leanpath{CollisionIdeals/Planar/GenericDegreeTwo/Main.lean}.

The two Palomar-facing declarations are
\leanpath{CollisionIdeals.Palomar.PaperTwo.explicitPlanarSecantProjector}
and
\leanpath{CollisionIdeals.Palomar.PaperTwo.quadraticPlanarCollisionRigidity}.
They are checked by
\texttt{lake build CollisionIdeals.Palomar.PaperTwo.Solution}.

The classical theorem in
Theorem~\ref{thm:keller-galois-rigidity} is neither installed as a Lean
axiom nor proved in the formalization.  Instead,
\leanpath{CollisionIdeals/General/Automorphism/GaloisRigidity.lean}
defines the proposition
\[
  \texttt{ComplexKellerGaloisRigidity 2}
\]
and passes a proof of it explicitly to every theorem using the literature
input.  Lean therefore proves the degree-two exclusion relative to this
single named hypothesis.

No normalization, ramification, boundary, purity, or finite-\'etale
interface occurs in the focused proof.  Those constructions belong to the
separate problem of nonnormal generic degree at least three.
\end{leanfiles}

\section{Conclusion}
\label{sec:conclusion}

For a planar Keller map, the divided-difference determinant gives
\[
  C_F\simeq A\times C_Fp_F,\qquad
  \Obs(F)=C_Fp_F,
\]
where \(p_F\) is the explicit off-diagonal collision projector.  Generic
degree two forces
\[
  K\otimes_BC_F\simeq_K L\times L,
\]
with diagonal evaluation equal to first projection, so \(p_F\neq0\).

The quadratic extension is nevertheless Galois.  Classical
Keller--Galois rigidity would make the map a polynomial automorphism,
forcing the collision scheme to equal its diagonal and \(p_F=0\).  This
incompatibility proves that no planar Keller map has generic degree two.

The remaining possible counterexamples to \(JC(2)\) have nonnormal
function-field extensions of generic degree at least three.  The degree-two
case itself is complete once the classical rigidity theorem is taken as
input.

\paragraph{Acknowledgment.}
OpenAI Codex (GPT-5.6 Sol) assisted with the Lean formalization and
manuscript.  The author reviewed and takes responsibility for all
statements, proofs, citations, and formalization code.

\bibliographystyle{alpha}
\bibliography{references}

\end{document}
